\documentclass[11pt,a4paper]{article}
\usepackage[utf8]{inputenc}
\usepackage{lmodern}
\usepackage[T1]{fontenc}
\usepackage[finnish]{babel}
\usepackage{graphicx}
\usepackage{epstopdf}
\usepackage{listings} 
\usepackage{xcolor}
\usepackage{textcomp}

\definecolor{dkgreen}{rgb}{0,0.6,0}
\definecolor{gray}{rgb}{0.5,0.5,0.5}
\definecolor{amber}{rgb}{1.0, 0.49, 0.0}

%Java lstset
\lstset{frame=none,
  language=Java,
  aboveskip=3mm,
  belowskip=3mm,
  showstringspaces=false,
  columns=flexible,
  basicstyle={\small\ttfamily},
  numbers=none,
  numberstyle=\tiny\color{gray},
  keywordstyle=\color{blue},
  commentstyle=\color{dkgreen},
  stringstyle=\color{amber},
  breaklines=true,
  breakatwhitespace=true,
  tabsize=3
}

%Multi-byte encoding
\lstset{literate=
  {á}{{\'a}}1 {é}{{\'e}}1 {í}{{\'i}}1 {ó}{{\'o}}1 {ú}{{\'u}}1
  {Á}{{\'A}}1 {É}{{\'E}}1 {Í}{{\'I}}1 {Ó}{{\'O}}1 {Ú}{{\'U}}1
  {à}{{\`a}}1 {è}{{\`e}}1 {ì}{{\`i}}1 {ò}{{\`o}}1 {ù}{{\`u}}1
  {À}{{\`A}}1 {È}{{\'E}}1 {Ì}{{\`I}}1 {Ò}{{\`O}}1 {Ù}{{\`U}}1
  {ä}{{\"a}}1 {ë}{{\"e}}1 {ï}{{\"i}}1 {ö}{{\"o}}1 {ü}{{\"u}}1
  {Ä}{{\"A}}1 {Ë}{{\"E}}1 {Ï}{{\"I}}1 {Ö}{{\"O}}1 {Ü}{{\"U}}1
  {â}{{\^a}}1 {ê}{{\^e}}1 {î}{{\^i}}1 {ô}{{\^o}}1 {û}{{\^u}}1
  {Â}{{\^A}}1 {Ê}{{\^E}}1 {Î}{{\^I}}1 {Ô}{{\^O}}1 {Û}{{\^U}}1
  {Ã}{{\~A}}1 {ã}{{\~a}}1 {Õ}{{\~O}}1 {õ}{{\~o}}1
  {œ}{{\oe}}1 {Œ}{{\OE}}1 {æ}{{\ae}}1 {Æ}{{\AE}}1 {ß}{{\ss}}1
  {ű}{{\H{u}}}1 {Ű}{{\H{U}}}1 {ő}{{\H{o}}}1 {Ő}{{\H{O}}}1
  {ç}{{\c c}}1 {Ç}{{\c C}}1 {ø}{{\o}}1 {å}{{\r a}}1 {Å}{{\r A}}1
  {€}{{\euro}}1 {£}{{\pounds}}1 {«}{{\guillemotleft}}1
  {»}{{\guillemotright}}1 {ñ}{{\~n}}1 {Ñ}{{\~N}}1 {¿}{{?`}}1
}


\begin{document}
\author{Ville Manninen
\thanks{mooc.fi: vw.manninen@gmail.com}
\thanks{opiskelijanumero: 014922187}}
\date{\today}
\title{Tietokantojen perusteet, kevät 2020 \\
Harjoitustyön raportti}
\maketitle
\begin{center}
\vspace{2cm}
\includegraphics{university-of-helsinki-2.eps} 
\end{center}
\newpage
\tableofcontents
\newpage
\section{Esittely}
Harjoitustyö on moniluokkainen komentoriviohjelma jonka avulla voidaan tallentaa asiakkaita, paikkoja, paketteja ja niihin liittyviä tapahtumia tietokantaan. Tietokantaan voidaan suorittaa ennakkoon määriteltyjä kyselyita, kyselyiden tulokset tulostetaan käyttäjän komentoriville. Ohjelmaan tallennetut tiedot jäävät muistiin tietokantaan ohjelman sulkemisen jälkeen.

\subsection*{Ohjelman ominaisuudet}
Ohjelmalla on kirjoitushetkellä seuraavat ominaisuudet.
\begin{itemize}
\item Ohjelman tietokannan luominen kun tietokanta puuttuu tai on puutteellinen.
\item Asiakkaiden lisääminen tietokantaan.
\item Paikkojen lisäämien tietokantaan.
\item Pakettien lisääminen tietokantaan.
\item Tapahtumien lisääminen tietokantaan.
\item Pakettien tapahtumien hakeminen.
\item Asiakkaan pakettien hakeminen ja yksittäisten pakettien tapahtumien lukumäärä.
\item Tapahtumien määrän hakeminen halutusta paikasta päivämäärällä.
\end{itemize}
\newpage
\subsection*{Komentorivi komennot}
Ohjelmalle annetaan syöttämällä komentoriville komentoa kuvaava numero(0-9).
\begin{description}
\item \textbf{1} lisää uuden asiakkaan tietokantaan.
\item \textbf{2} lisää uuden paikan tietokantaan.
\item \textbf{3} lisää uuden paketin tietokantaan.
\item \textbf{4} lisää uuden tapahtuman tietokantaan.
\item \textbf{5} tulostaa paketin tapahtumat tietokannasta.
\item \textbf{6} tulostaa asiakkaan paketit ja niihin liittyvien tapahtumien määrän tietokannasta.
\item \textbf{7} tulostaa tapahtumien määrän haku paikasta tietyllä päivämäärällä tietokannasta.
\item \textbf{8} suorittaa tehokkuustestin ja tulostaa testin tulokset käyttäjälle.
\item \textbf{9} luo tietokannan.
\item \textbf{0} Sulkee ohjelman.
\end{description}

\newpage
\section{Kaaviot ja Skeemat}
\subsection*{Tietokantakaavio}

\subsection*{SQL-skeema}

\newpage
\section{Tehokkuustesti}
Ohjelmaan sisältyy automaattinen tehokkuustesti jonka avulla voidaan testata tietokannan tehokkuutta. Tehokkuustesti suorittaa ennalta määriteltyjä muutoksia tietokantaan ja ilmoittaa niihin käytetyn ajan. Testi koostuu kuudesta eri vaiheesta. 

Testin vaiheissa yhdestä neljään tietokantaan lisätään tuhat asiakasta, paikkaa ja pakettia. Seuraavaksi tietokantaan lisätään miljoona tapahtumaa käyttäen satunnaisesti tietokantaan lisättyjä arvoja. Viimeisissä vaiheissa viisi ja kuusi testi suorittaa kyselyitä joissa haetaan tuhannen asiakkaan pakettien määrä, sekä tuhannen paketin tapahtumien määrä.

\subsection*{Ilman indeksejä}
Ohjelmaa testatessa ilman ohjelmoituja hakemistorakenteita tuhannen arvon lisääminen tapahtui alle sekunnissa, ja miljoonaa arvoa lisätessä aikaa kului 10-12 sekunttia. 

Kyselyitä suorittaessa tehokkuuteen vaikutti merkittävästi haettavan taulukon koko. Haettaessa asiakkaan pakettien määrää tuhannen arvon taulukosta aikaa kului alle sekuntti, Kun haettiin paketin tapahtumien määrää miljoonan arvon taulukosta aikaa kului noin minuutti.


\subsection*{Indeksien lisäämisen jälkeen}
Seuraavaksi ohjelman tehokkuustestiin lisättiin hakemistorakenne jonka avulla tapahtumat taulukolle luotiin indksi sarakkeen arvoille 'paketti.id'. Lisäyksen jälkeen tapahtumien haku paketin arvolla toimii huomattavasti tehokkaammin, lisäämisen jälkeen paketin tapahtumien määrän hakeminen, johon ennen kului aikaa yli minuutin, tapahtui alle sekunnissa.

Huomattavaa oli myös se missä kohdassa testiä indeksi lisättiin, kun indeksi lisättiin ennen miljoonan tapahtuman lisäämistä mainittujen tapahtumien luomiseen kului tuplasti enemmän aikaa. 

\newpage
\section{Tietokannan eheys}
Ohjelman tietokannan eheyttä pidetään yllä käyttämällä SQLiten sekä Javan tarjoamia ominaisuuksia.
Tietokantaan ei ole mahdollista luoda samannimisiä asiakkaita tai paikkoja, eikä tietokantaan luoduille paketeille saa antaa samaa seurantakoodia. Samannimisiä syötteitä valvotaan SQLiten rajoitteella UNIQUE, joka varmistaa ettei samaan tauluun lisätä samannimisiä arvoja. Lisäksi SQLiten rajoitteella NOT NULL valvotaan että sarakkeisiin ei lisätä tyhjiä arvoja.
 
Java-ohjelma tarkistaa käyttäjän antamat syötteet ja ilmoittaa syyn jos annettu syöte on puutteellinen, esimerkiksi jos käyttäjän syöte on tyhjä. Tietokantaan lisättäessä paketteja tai tapahtumia jotka viittaavat muualle tallennettuun tietoon, sovellus varmistaa että annetut viitteet löytyvät tietokannasta.

\newpage
\section{Lähdekoodi}
%Tarkista lstinputlisting polut
\textbf{https://github.com/Viltska/sql-app-training}.

\subsection*{Main.java}
\lstinputlisting[language=Java]{E:/Koulu/Tietokantojen Perusteet/sql-app-training/src/Main.java}
\subsection*{DatabaseManager.java}
\lstinputlisting[language=Java]{E:/Koulu/Tietokantojen Perusteet/sql-app-training/src/DatabaseManager.java}
\subsection*{Asiakkaat.java}
\lstinputlisting[language=Java]{E:/Koulu/Tietokantojen Perusteet/sql-app-training/src/Asiakkaat.java}
\subsection*{Paikat.java}
\lstinputlisting[language=Java]{E:/Koulu/Tietokantojen Perusteet/sql-app-training/src/Paikat.java}
\subsection*{Paketit.java}
\lstinputlisting[language=Java]{E:/Koulu/Tietokantojen Perusteet/sql-app-training/src/Paketit.java}
\subsection*{Tapahtumat.java}
\lstinputlisting[language=Java]{E:/Koulu/Tietokantojen Perusteet/sql-app-training/src/Tapahtumat.java}
\end{document}